% ============================================================
% Notes: Random Walks, Normal/Anomalous Diffusion, Lévy Walks
% + link to fractional calculus / nonlocality
% Sources:
%  - Herrmann, Fractional Calculus: An Introduction for Physicists
%    (motivation: anomalous diffusion + fractional Brownian motion) :contentReference[oaicite:1]{index=1}
%  - Standard stochastic process theory (Brownian motion, CTRW, Lévy walks)
% ============================================================

\section{Random walks and diffusion processes}

\subsection{Discrete-time random walk (1D)}
Let $X_n$ be the position after $n$ steps:
\begin{equation}
X_n = \sum_{k=1}^n \xi_k,
\qquad 
\xi_k \in \{+a,-a\},
\qquad 
\mathbb{P}(\xi_k=+a)=p,\ \mathbb{P}(\xi_k=-a)=q=1-p.
\end{equation}

Mean and variance:
\begin{equation}
\mathbb{E}[X_n] = n a (p-q),
\qquad
\mathrm{Var}(X_n)=4npq\,a^2.
\end{equation}

Symmetric walk ($p=q=1/2$):
\begin{equation}
\mathbb{E}[X_n]=0,
\qquad
\mathrm{Var}(X_n)=na^2.
\end{equation}

\subsection{Central limit theorem (Gaussian scaling)}
If $\xi_k$ are i.i.d. with finite variance $\sigma^2$, then for large $n$:
\begin{equation}
\frac{X_n - n\mu}{\sqrt{n\sigma^2}}
\ \xrightarrow[n\to\infty]{d}\
\mathcal{N}(0,1),
\qquad
\mu=\mathbb{E}[\xi_k].
\end{equation}

So the displacement distribution approaches a Gaussian.

% ------------------------------------------------------------

\section{Normal diffusion}

\subsection{Mean squared displacement (MSD)}
For a stochastic process $X(t)$ define
\begin{equation}
\langle X^2(t)\rangle := \mathbb{E}[X(t)^2].
\end{equation}

\textbf{Normal diffusion} is characterized by linear growth:
\begin{equation}
\langle X^2(t)\rangle \sim 2Dt \quad \text{(1D)}, 
\qquad
\langle r^2(t)\rangle \sim 2dDt \quad \text{(in $d$ dimensions)}.
\end{equation}

\subsection{Diffusion equation (Fick / heat equation)}
Let $P(x,t)$ be the probability density of being at $x$ at time $t$.
Normal diffusion satisfies:
\begin{equation}
\frac{\partial P}{\partial t} = D\frac{\partial^2 P}{\partial x^2}.
\end{equation}

With initial condition $P(x,0)=\delta(x)$, the Green function is Gaussian:
\begin{equation}
P(x,t)=\frac{1}{\sqrt{4\pi Dt}}\exp\!\left(-\frac{x^2}{4Dt}\right).
\end{equation}

% ------------------------------------------------------------

\section{Anomalous diffusion}

\subsection{Definition via MSD scaling}
\textbf{Anomalous diffusion} is defined by a power-law MSD:
\begin{equation}
\langle X^2(t)\rangle \sim t^\alpha,
\qquad \alpha\neq 1.
\end{equation}

Classification:
\begin{itemize}
\item $0<\alpha<1$: \textbf{subdiffusion} (slower than normal).
\item $\alpha=1$: \textbf{normal diffusion}.
\item $1<\alpha<2$: \textbf{superdiffusion}.
\item $\alpha=2$: \textbf{ballistic transport} ($X(t)\sim vt$).
\end{itemize}

\subsection{Physical origins}
Common mechanisms:
\begin{itemize}
\item \textbf{Trapping / waiting times} (long rests) $\Rightarrow$ subdiffusion.
\item \textbf{Long jumps / flights} (rare huge displacements) $\Rightarrow$ superdiffusion.
\item \textbf{Long-time correlations} (persistent velocity) $\Rightarrow$ superdiffusion.
\item \textbf{Fractal / disordered media} $\Rightarrow$ anomalous scaling.
\end{itemize}

\subsection{Connection to fractional calculus (Herrmann motivation)}
Herrmann highlights that anomalous diffusion and fractional Brownian motion
became a major physics motivation for fractional methods (especially since the 1980s). :contentReference[oaicite:2]{index=2}

Core idea:
\begin{itemize}
\item \textbf{Temporal nonlocality / memory} $\Rightarrow$ fractional \emph{time} derivatives.
\item \textbf{Spatial nonlocality / long jumps} $\Rightarrow$ fractional \emph{space} derivatives.
\end{itemize}

% ------------------------------------------------------------

\section{Continuous-time random walk (CTRW) framework}

\subsection{CTRW definition}
A CTRW consists of:
\begin{itemize}
\item jump lengths $\eta_1,\eta_2,\dots$ with PDF $\lambda(\eta)$
\item waiting times $\tau_1,\tau_2,\dots$ with PDF $\psi(\tau)$
\end{itemize}
and position after $N(t)$ jumps:
\begin{equation}
X(t)=\sum_{k=1}^{N(t)}\eta_k,
\qquad
N(t)=\max\left\{n:\sum_{k=1}^n \tau_k \le t\right\}.
\end{equation}

\subsection{Normal diffusion condition}
If $\mathbb{E}[\tau]<\infty$ and $\mathbb{E}[\eta^2]<\infty$,
then at large times:
\[
\langle X^2(t)\rangle \sim t.
\]

\subsection{Subdiffusion from heavy-tailed waiting times}
If the waiting time distribution has a power-law tail:
\begin{equation}
\psi(\tau)\sim \tau^{-1-\beta},
\qquad 0<\beta<1,
\end{equation}
then $\mathbb{E}[\tau]=\infty$ and typically
\begin{equation}
\langle X^2(t)\rangle \sim t^\beta
\qquad (\text{subdiffusion}).
\end{equation}

This is a standard route to fractional-in-time diffusion equations.

\subsection{Superdiffusion from heavy-tailed jump lengths}
If jump lengths have heavy tails:
\begin{equation}
\lambda(\eta)\sim |\eta|^{-1-\mu},
\qquad 0<\mu<2,
\end{equation}
then $\mathbb{E}[\eta^2]=\infty$ and the process can become superdiffusive.
This is associated with Lévy stable statistics and fractional Laplacians.

% ------------------------------------------------------------

\section{Lévy flights vs Lévy walks}

\subsection{Lévy flight (mathematical idealization)}
A \textbf{Lévy flight} is a random walk with i.i.d.\ heavy-tailed step lengths:
\begin{equation}
\lambda(\eta)\sim |\eta|^{-1-\mu},\qquad 0<\mu<2.
\end{equation}

Key property: the variance diverges:
\begin{equation}
\mathbb{E}[\eta^2]=\infty.
\end{equation}

Problem: steps occur in effectively zero time, so arbitrarily long jumps happen instantly
(unphysical for finite-speed transport).

\subsection{Lévy walk (physically meaningful)}
A \textbf{Lévy walk} couples displacement and travel time.
A typical model:
\begin{equation}
\eta = v\,\tau \cdot s,
\qquad s\in\{+1,-1\},
\end{equation}
where $\tau$ is a flight time and $v$ is a finite speed.

If flight times are heavy-tailed:
\begin{equation}
\psi(\tau)\sim \tau^{-1-\gamma},
\end{equation}
then transport can be superdiffusive but still respects finite propagation speed.

\subsection{Scaling behavior (qualitative)}
Depending on $\gamma$:
\begin{itemize}
\item large $\gamma$ (short flights dominate) $\Rightarrow$ normal diffusion
\item intermediate $\gamma$ $\Rightarrow$ superdiffusion
\item very small $\gamma$ $\Rightarrow$ nearly ballistic scaling
\end{itemize}

\subsection{Why Lévy walks matter}
Lévy walks provide:
\begin{itemize}
\item heavy tails (rare long moves)
\item finite velocity constraint
\item realistic superdiffusion model in physics/biology
\end{itemize}

% ------------------------------------------------------------

\section{Fractional diffusion equations (bridge to anomalous diffusion)}

\subsection{Fractional time diffusion (subdiffusion)}
Subdiffusion with memory can be modeled by a fractional time derivative:
\begin{equation}
\prescript{C}{}D_t^\beta P(x,t)
=
D\frac{\partial^2 P}{\partial x^2},
\qquad 0<\beta<1.
\end{equation}

This encodes \textbf{nonlocality in time} (history dependence).

\subsection{Fractional space diffusion (superdiffusion)}
Superdiffusion with long jumps can be modeled using a fractional Laplacian:
\begin{equation}
\frac{\partial P}{\partial t}
=
- K_\mu\,(-\Delta)^{\mu/2}P,
\qquad 0<\mu<2.
\end{equation}

This encodes \textbf{nonlocality in space} (long-range jumps).

\subsection{Combined space-time fractional transport}
More general anomalous transport can involve both:
\begin{equation}
\prescript{C}{}D_t^\beta P(x,t)
=
- K_\mu\,(-\Delta)^{\mu/2}P.
\end{equation}

% ------------------------------------------------------------

\section{High-yield summary}
\begin{itemize}
\item \textbf{Random walk} $\Rightarrow$ diffusion in the continuum limit.
\item \textbf{Normal diffusion:} $\langle X^2(t)\rangle \sim t$ and Gaussian $P(x,t)$.
\item \textbf{Anomalous diffusion:} $\langle X^2(t)\rangle \sim t^\alpha$, $\alpha\neq 1$.
\item \textbf{Subdiffusion} comes from long waiting times / trapping $\Rightarrow$ memory.
\item \textbf{Superdiffusion} comes from long jumps / persistent motion.
\item \textbf{Lévy flight:} heavy-tailed jumps, infinite variance, infinite-speed issue.
\item \textbf{Lévy walk:} heavy-tailed \emph{flight times} with finite velocity (physical).
\item \textbf{Fractional calculus link (Herrmann):} anomalous diffusion motivates fractional operators
as nonlocal descriptions of transport. :contentReference[oaicite:3]{index=3}
\end{itemize}
